\chapter{Related Work}
\label{chapter:relacionados}

This chapter positions \textit{G-FShield} in the literature on \acp{ids} for \acp{cps}, GRASP-FS, distributed architectures, and observability.

\section{Intrusion Detection in Cyber-Physical Systems}

\citeonline{Cardenas2009}, \citeonline{Mitchell2014}, and \citeonline{Alcaraz2019} show that CPS security must consider the coupling of software, networks, and physical assets. In digital substations, the \ac{iec} literature discusses specific traffic and timing requirements \cite{iec61850,mackiewicz2006overview,aftab2020iec}. \citeonline{quincozes2021survey} and \citeonline{quincozes2022ereno} stress the importance of realistic corpora for IDS evaluation.

\section{Feature Selection and GRASP-FS}

\citeonline{Hall1999,Saeys2007,guyon2003introduction} present feature selection as a technique for reducing redundancy and processing cost. Outside the specific lineage of this dissertation, \citeonline{yusta2009different} compares metaheuristic strategies for feature selection, and \citeonline{esseghir2010effective} combines filters and wrappers in a GRASP scheme. These works address search design but not operation of an observable distributed pipeline.

GRASP is grounded by \citeonline{ResendeRibeiro2010}; its use for IDS feature selection is investigated by \citeonline{quincozes2020grasp,quincozes2021performance,quincozes2022extended}. Recent surveys indicate that metaheuristics remain relevant to attack detection in industrial and IIoT environments \cite{Termanini2025Survey,Francis2026MFSReview}. This trajectory supports the algorithmic components reused by G-FShield, not a claim of a new algorithm.

As an external comparator, \citeonline{Scherer2024IWSHAP} propose IWSHAP, which ranks features with SHAP values and applies incremental wrapper selection with XGBoost to CAN traffic. For its X-CANIDS suspension experiment, the paper reports positive-class F1 of approximately 0.9186 with 16 of the 688 features. These are published results obtained under the work's own protocol, classifier, and data scale; in this dissertation, they therefore contextualize selection quality but are not directly subtracted from G-FShield metrics. The controlled experimental comparison uses the public artifact pinned to a revision, common train, validation, and test partitions, and the same final evaluator for every subset.

\section{Distributed Architectures and Observability}

Earlier work explored parallelism in parts of GRASP-FS. \citeonline{Silva2023GDLSFS} distributes local search, whereas \citeonline{Naves2023DRG} distributes restricted-candidate-list generation. In metaheuristics from another domain, \citeonline{SUBRAMANIAN20101899} exemplifies parallel search but does not address feature selection or the operational integration studied here.

In another line of work, distributed-systems research grounds component decoupling \cite{Tanenbaum2007Distributed}. Service decomposition is discussed by \citeonline{Newman2021}, and messaging logs by \citeonline{Kreps2011}. Site-reliability work addresses operational telemetry \cite{beyer2016site,mateus2021security}. Observability surveys discuss diagnosis in service architectures \cite{usman2022observabilitysurvey,Faseeha2025MicroservicesObs}. \citeonline{wu2012online}, in turn, addresses feature selection from streaming features. That is a different problem: the evaluated G-FShield receives a pre-materialized corpus and does not perform online selection. Among the works examined, none combines DRG, DLS, event persistence, and operational queries in a GRASP-FS flow; because the search was not systematic, this conclusion is limited to the set actually analyzed.

\section{Identified Gap}

\textit{G-FShield} integrates distributed generation, distributed local search, result persistence, and telemetry visualization in one architecture. Its central contribution is architectural and software-engineering-oriented: end-to-end integration, orchestration, and process traceability. DRG and DLS build on prior work and are not presented as a new search algorithm.

\section{Comparative Table}

Table~\ref{tab:related_work_gfshield} summarizes the discussed capabilities. ``Distributed construction'' means distributing initial-solution formation; ``distributed refinement'', distributing local search; ``E2E integration'', connecting both in one flow; and ``telemetry'', persisting signals that can reconstruct an execution. The table does not attribute online selection to G-FShield.

\begin{landscape}
\begin{table}[p]
\centering
\scriptsize
\caption{Summary of related work cited in this manuscript.}
\label{tab:related_work_gfshield}
\setlength{\tabcolsep}{4pt}
\renewcommand{\arraystretch}{1.15}
\begin{tabular}{>{\raggedright\arraybackslash}p{.20\linewidth}*{6}{>{\centering\arraybackslash}p{.105\linewidth}}}
\toprule
\textbf{Reference} & \textbf{CPS/ IDS} & \textbf{GRASP- FS} & \shortstack{\textbf{Distributed}\\\textbf{construction}} & \shortstack{\textbf{Distributed}\\\textbf{refinement}} & \shortstack{\textbf{E2E}\\\textbf{integration}} & \textbf{Telemetry} \\
\midrule
Resende and Ribeiro (2016) & -- & \checkmark & -- & -- & -- & -- \\
Yusta (2009); Esseghir et al. (2010) & -- & Partial & -- & -- & -- & -- \\
Quincozes et al. (2020--2022) & \checkmark & \checkmark & -- & -- & -- & -- \\
Scherer et al. (2024), IWSHAP & \checkmark & -- & -- & -- & -- & -- \\
Silva et al. (2023) & \checkmark & \checkmark & -- & \checkmark & Partial & -- \\
Faria et al. (2023) & \checkmark & \checkmark & \checkmark & -- & Partial & -- \\
Beyer et al. (2016); Usman et al. (2022) & -- & -- & -- & -- & -- & \checkmark \\
Wu et al. (2013) & -- & -- & -- & -- & -- & -- \\
\textbf{This work (G-FShield)} & \checkmark & \checkmark & \checkmark & \checkmark & \checkmark & \checkmark \\
\bottomrule
\end{tabular}
\end{table}
\end{landscape}

The table therefore supports an architectural distinction: G-FShield connects feature-selection stages and their operational signals. A claim of exhaustive literature coverage would require a systematic review or a documented bibliographic strategy.
